\documentclass[noanswers, addpoints, 12pt]{exam}
\usepackage[T1]{fontenc}
\usepackage[utf8]{inputenc}
\usepackage{lmodern}
\usepackage{tikz}
\usetikzlibrary{positioning,shapes.multipart,arrows.meta,shadows}
\usepackage[french]{babel}
\usepackage{xcolor}
\usepackage{listings}
\usepackage{graphicx}
\usepackage[most]{tcolorbox}

\newcommand\emptybox[1]{%
 \hbox{\smash{\hspace*{\dimexpr-\hoffset}%
  }\framebox(\textwidth, #1){}}}

\definecolor{codegreen}{rgb}{0,0.6,0}
\definecolor{codegray}{rgb}{0.5,0.5,0.5}
\definecolor{codepurple}{rgb}{0.58,0,0.82}
\definecolor{backcolour}{rgb}{0.95,0.95,0.92}
 
\lstdefinestyle{mystyle}{
    backgroundcolor=\color{backcolour},   
    commentstyle=\color{codegreen},
    keywordstyle=\color{magenta},
    numberstyle=\tiny\color{codegray},
    stringstyle=\color{codepurple},
    basicstyle=\footnotesize,
    breakatwhitespace=false,         
    breaklines=true,                 
    captionpos=b,                    
    keepspaces=true,                 
    numbers=left,                    
    numbersep=5pt,                  
    showspaces=false,                
    showstringspaces=false,
    showtabs=false,                  
    tabsize=2
}
\lstset{style=mystyle, language=Python}

\title{CC 1er NSI - Fonction}
\author{Babash - NSI Tani Malandi et Johanne Bakalara}
\date{Novembre 2024, Durée: 1h30min}


\lstset{extendedchars=true,literate={à}{{\`a}}1 {é}{{\'e}}1}
\begin{document}
\runningfooter{}{\footnotesize Babash - NSI Tani Malandi et Johanne Bakalara \textperiodcentered{} CC BY-SA 4.0 \textperiodcentered{} page \thepage}{}
\firstpagefooter{}{\footnotesize Babash - NSI Tani Malandi et Johanne Bakalara \textperiodcentered{} CC BY-SA 4.0 \textperiodcentered{} page \thepage}{}

\maketitle

\pagestyle{headandfoot}
\thispagestyle{headandfoot}
\firstpageheader{Nom : \\  Prénom :\\ Classe :}{}{}


\section*{Consignes:}
\begin{center} 
\fbox{\parbox{\linewidth}{
\begin{center}
Le devoir est à rédiger sur le sujet.\\
La clarté de la rédaction et les explications font partie intégrante de la notation. Une réponse non rédigée sans explication sera considérée fausse.\\
\end{center}


}}
\end{center}

\vsword{Note:}
\hsword{Note:}

\section*{Barème :}
\begin{center}
  \par\bigskip
  \gradetable[h]
\end{center}
\begin{questions}

\vspace{0.5cm}
\section{Question de cours : }

\question[1] Une fonction est-elle une structure de contrôle : Vrai ou Faux ?
\question[3] Soit le code ci-dessous :
\begin{lstlisting}[style=mystyle, language=Python]
def somme_entre_deux_nombres(a, b):
    somme = 0
    for i in range(a, b):
        somme += i
    return somme
\end{lstlisting}
\begin{parts}
    \part Quels sont les paramètres de la fonction ci-dessus ?
    \part Quels sont les variables locales de la fonction ci-dessus ?
    \part Quelles sont les valeurs de retour de la fonction ci-dessus ?
\end{parts}
\question[2]Écrire un appel de la fonction \textit{somme\_entre\_deux\_nombres} pour a=10 et b=15. On prendra soin de stocker les éventuelles valeurs de retour dans des variables.
\question[4] Associer à chacun des chiffres du schéma ci-dessous un des termes suivant : appeler, arguments, valeur de retour, corps, en-tête, définir, paramètre, nom. \begin{figure}[h]
    \centering
    \begin{center}
\begin{tikzpicture}[font=\small,
  num/.style={draw,minimum width=0.7cm,minimum height=0.55cm,inner sep=1pt},
  lab/.style={anchor=west,inner sep=2pt},
  box/.style={draw=cyan!60!blue,rounded corners=6pt,very thick,inner sep=7pt}]
\node[box,anchor=center,align=left,font=\ttfamily\small] (g) at (0,0) {%
\textcolor{blue!70!black}{def} nbrCaractere (ch):\\
\hspace*{1.4em}i=0\\
\hspace*{1.4em}\textcolor{orange!80!black}{for} j \textcolor{orange!80!black}{in} ch:\\
\hspace*{2.8em}i+=1\\
\hspace*{1.4em}\textcolor{brown!80!black}{return} (i)};
\node[box,font=\ttfamily\small] (d) at (8,0) {nbrCaractere ('Alex')};
\node[num] (n1) at (-3,3.3) {1}; \node[lab] at (n1.east) {une fonction};
\node[num] (n2) at (5.6,1.6) {2}; \node[lab] at (n2.east) {une fonction};
\node[num] (n6) at (-1.6,2.3) {6}; \node[lab] at (n6.east) {de la fonction};
\node[num] (n8) at (2.6,2.3) {8};
\node[num] (n3) at (-6.6,1.0) {3}; \node[lab] (t3) at (n3.east) {de la fonction};
\node[num] (n4) at (-6.6,0) {4}; \node[lab] (t4) at (n4.east) {de la fonction};
\node[num] (n5) at (-6.6,-1.0) {5};
\node[num] (n7) at (8.3,-1.8) {7}; \node[lab] at (n7.east) {de la fonction};
\draw[->,thick] (t3.east) -- (-2.5,1.0);
\draw[->,thick] (t4.east) -- (-2.5,0);
\draw[->,thick] (n5.east) -- (-2.5,-1.0);
\draw[->,thick] (n6.south) -- (-0.6,1.15);
\draw[->,thick] (n8.south) -- (1.5,1.15);
\draw[->,thick] (n7.north) -- (8.3,-0.45);
\end{tikzpicture}
\end{center}

    \label{fig:schema-a-trou}
\end{figure}

\section{Compréhension de code : }
\begin{lstlisting}[style=mystyle, language=Python]
import random

def demander_nombre():
    rep = int(input("quelle est votre choix ?: "))
    return rep

def aide(rep, nbHasard):
    if rep<nbHasard:
        print("+")
    if rep>nbHasard:
        print("-")
        
essai = 11
#on crée une variable qui va stocker le nombre à trouver .
nbHasard = random.randint(1 , 100)
rep = 0

while rep != nbHasard and essai > 0:
    essai = essai - 1
    rep = demander_nombre()
    aide(rep, nbHasard)

if essai > 0 :
    print("Bravo, vous avez trouvé la bonne réponse !")
else :
    print("Vous avez perdu !")
\end{lstlisting}
\question[3] Dans le code ci-dessus, entourer les définitions des fonctions (en entière) et souligner les appels.

\begin{lstlisting}[style=mystyle, language=Python]
from turtle import *

def surprise(cote, repetition):
  # Au depart la tortue est a l'horizontale.
  for i in range(repetition):
    forward((i//4+1)*cote) # forward sert a avancer
    right(90) # right sert a tourner a droite

surprise(50, 16)
\end{lstlisting}
\question[2] Dessiner ce qu'affiche le code ci-dessus. 

\begin{lstlisting}[style=mystyle, language=Python]
def f(x, y, z):
    return x + y, x + z, y + z

def g(x, y):
    if x > 10 :
        return x, y
    else :
        return y, x

def combo(x, y, z):
    s1, s2, s3  = f(x, y, z)
    print(s1, s2, s3)
    a, b = g(s1, s3)
    print(a, b)
    return s1 - a, s2 - b 

res1, res2 = combo(4 ,9 ,2)
print(res1, res2)
\end{lstlisting}
\question[3] Qu'affiche le code ci-dessus ?\\

\section{Rédaction de fonction :}
\question[2] Écrire une fonction division qui renvoie le quotient de deux nombres donnés en paramètres, attention, vous devrez prendre garde que le diviseur soit différent de 0.
\end{questions}
\end{document}

